% !TEX root = ../5.DOUBLE-PIPE.tex
\section{DISCUSSION}\label{discussion}

\subsection*{Question 1: Is heat loss significant? Why?}\label{q1-heat-loss}

From calculation, it can be seen that the heat loss is significant, more detail, the heat loss in tube B is generally more than the one in tube C.

\textbf{Causes:}
\begin{itemize}
    \item Heat loss is the result of heat transfer from the fluids to the environment, it may be because of a lack of good insulation.
    \item The heat is not only loss to the environment but also is 'trapped' in the metal wall. Also, copper has very high thermal-conductivity.
    \item The flow rates are not really stable during the experiment, due to the instability in pump work, which leads to unstable heat transfer.
    \item Tube B has more heat loss than tube C because the fluid streams in tube B contact each other in vertical flow leading to shorter heat transfer time in comparison with tube C. Beside that, fluid leaking on tube B is also a cause.
    \item Temperature values were displayed and observed subjectively leading to some cases that the inlet and outlet temperature were equal, which means there was no heat transfer, which is incorrect and does not reflect the nature of the heat exchanging process.
\end{itemize}

\subsection*{Question 2: Errors percentage and their causes? How to avoid them?}\label{q2-errors}

Errors percentage is significant. There are several causes:
\begin{itemize}
    \item Flow rates of the hot and cold streams are not stable during the experiment. It is impossible to maintain a constant flow rate.
    \item The temperature changes continuously leading to errors in observing results.
    \item The hot stream is refluxed continuously leading to a non-constant hot stream inlet temperature.
    \item The higher the temperature, the greater the heat loss.
    \item Insulation is not good enough.
    \item There is fluid leaking during the experiment.
    \item Errors occur in flow meter and thermal-meter (system error).
    \item The experiment is performed once only, as the result, the repetition of the results is not confirmed.
    \item Students made mistakes in performing the experiment.
    \item Rounding number also contributes to the errors. The observed results are average numbers, which may leads to inaccurate calculation.
\end{itemize}

\textbf{How to avoid errors:}
\begin{itemize}
    \item Supply a good insulation cover, repair leaks on the tubes.
    \item Perform the experiment in not-too-high temperature.
    \item Observe and note down the results quickly.
    \item Repeat the experiment more than once.
    \item Remove fouling inside tubes if possible.
\end{itemize}

\subsection*{Question 3: Compare experimental heat transfer coefficient $K_l$ with overall heat transfer coefficient $K_l^*$.}\label{q3-comparison}

\begin{itemize}
    \item In this experiment, $K_l$ is greater than $K_l^*$ in most cases, which may be impossible. Because errors during the experiment are much greater than heat loss.
    \item The average wall temperature is estimated during calculation leading to incorrect $Pr_{wall}$ in comparison with the experiment.
    \item The complex heat transfer had the student ignore some factor and only note down the main factors. For example, convectional heat transfer in tube B was ignored despite of the fact that convectional heat transfer plays an important role in how heat transfer is functioned.
    \item The difference between $K_l$ and $K_l^*$ is the result of significant heat loss, heat transfer from outside, fouling, \dots During calculations, fouling was ignored since it is impossible to measure the fouling, which leads to noticeable differences between experimental and overall K.
    \item While $K_l^*$ keeps its stability and slightly increases linearly, $K_l$ fluctuates in the most cases.
    \item The experimental heat transfer coefficient of cold streams is lower than the overall due to the low heat transfer, because heat transfer area is small in comparison with flow rate, and there may be also fouling.
    \item The experiment heat transfer coefficient of hot stream is higher than the overall due to the heat is absorbed by the copper wall, more over, copper has higher thermal-conductivity than water.
\end{itemize}

\textbf{In conclusion:} The high thermal-conductivity copper wall which is covered by fouling is an important reason for heat loss, heat 'trapped' leading to huge difference between overall heat transfer coefficient and experimental one.
